\section{结论}

\subsection{主要结论}

本文以Quanser五杆绘图机器人的实际拓扑和笔式绘迹场景为参照，建立了竖直板面椭圆轮廓描绘的课程设计模型。通过“解除闭环分析2R支链，再恢复闭环分析整机”的路线，完成了串联正逆运动学、雅可比、拉格朗日与牛顿—欧拉动力学，以及五杆并联正逆运动学和奇异性分析，并将支链动力学投影到整机的两个主动关节。

两种独立串联动力学实现、正动力学回放和能量关系在给定理想模型内保持一致，说明本例的符号、惯性耦合和载荷计数具有可重复的计算依据。对五杆机构而言，两条支链同时可达只是必要的几何条件；并联侧约束矩阵仍可能退化，使机构在关节锁定时出现瞬时自由运动，并使某些方向的传力性能恶化。本例相同$1\units N$竖直等效力对照进一步隔离了几何作用，说明不能仅靠逆解存在判断轮廓是否适宜执行。

与实际描绘任务联系起来，机构安装方向决定面内重力负载，执行时间决定惯性和速度二次项的量级，而近奇异位形的静态传力困难不能通过单纯降速消除。因此，轮廓位置、模式连续性、奇异性裕量与驱动力矩应共同参与轨迹设计。

\subsection{后续工作}

下一步可在真实绘图平台上测量几何参数和输出端摩擦，加入电机及减速器模型，比较预测力矩与实际电流换算力矩；随后引入笔尖接触摩擦及压笔机构，检验不同描绘速度下的轨迹偏差。这些实验将回答当前理想刚体模型未覆盖的参数误差与接触效应，同时保留本文的运动学与受约束动力学框架作为计算基础。

\subsection{代码与数据获取}

本文配套的Python仿真代码、计算数据、配图、LaTeX源文件与报告PDF已整理至作者GitHub账号Zhangzianqi的项目仓库：
\begin{center}
\small\url{https://github.com/Zhangzianqi/robot-dynamics-fivebar}
\end{center}
仓库README说明了文件结构、运行依赖及计算顺序，可用于重算本文轨迹、动力学与奇异性分析结果。
